% =============================================================================
% main.tex — single-paper deep-dive Beamer deck (Simple theme)
%
% Copy this file into:  ./slides/<paper-short-name>/main.tex
% Compile with XeLaTeX (preferred):  make   (or xelatex twice)
%
% Path assumptions (see the Makefile in this template):
%   - This deck lives at      <workspace>/slides/<paper-short-name>/
%   - The theme lives at the workspace root:  <workspace>/beamerthemeSimple.sty
%   - The Makefile sets TEXINPUTS=../../: so xelatex finds the theme .sty.
%   - \graphicspath below makes \includegraphics find this deck's local figures/
%     (and the workspace root, for any shared images you keep there).
%
% Theme: this deck loads the workspace-wide theme from <workspace>/theme.tex if it
% exists, otherwise falls back to the bundled Simple theme. To restyle ALL decks at
% once, edit <workspace>/theme.tex (see the workspace AGENTS.md "Theming"). To restyle
% only THIS deck, replace the \InputIfFileExists line below with your own \usetheme{...}.
% =============================================================================
\documentclass[aspectratio=169]{beamer}
\InputIfFileExists{../../theme.tex}{}{\usetheme{Simple}}

% Resolve this deck's local figures/ (and the workspace root for shared images).
\graphicspath{{figures/}{../../}}

\usepackage{booktabs}   % nicer tables
\usepackage{amsmath}

% ---------- talk metadata (fill in) ----------
\title{<Paper Title>}
\subtitle{<One-sentence takeaway>}
\author{Presenter: <your name>}
\institute{Paper: <Authors>, <Venue/Year>}
\date{\today}

\begin{document}

% ---------- 1. Title ----------
\begin{frame}[plain]
  \titlepage
\end{frame}

% ---------- Outline ----------
\begin{frame}{Outline}
  \tableofcontents
\end{frame}

% =============================================================================
% Fill in the sections below. Keep ONE main idea per slide. Split crowded
% slides instead of shrinking fonts. Compile + visually inspect after each
% section (see SKILL.md sections 6 and 10).
% =============================================================================

% ---------- 2. Motivation ----------
\section{Motivation}
\begin{frame}{Motivation: The Problem}
  \begin{itemize}
    \item What problem does the paper address?
    \item Why does it matter?
    \item Why was existing work insufficient?
  \end{itemize}
\end{frame}

% ---------- 3. Background ----------
\section{Background}
\begin{frame}{Background}
  % Enough context for a grad student who has NOT read the paper.
  \begin{itemize}
    \item Key concept 1
    \item Key concept 2
  \end{itemize}
\end{frame}

% ---------- 4. Problem Formulation ----------
\section{Problem Formulation}
\begin{frame}{Problem Formulation}
  \begin{itemize}
    \item Target system / setting
    \item Assumptions and constraints
    \item What is being optimized
  \end{itemize}
\end{frame}

% ---------- 5. Key Insight ----------
\section{Key Insight}
\begin{frame}{Key Insight}
  % Prefer a diagram here over text.
  \begin{itemize}
    \item What did the authors notice?
    \item Why is it non-obvious?
    \item Why does the approach make sense?
  \end{itemize}
\end{frame}

% ---------- 6. Design / Method ----------
\section{Design / Method}
\begin{frame}{Design Overview}
  % Split complex mechanisms across multiple slides (Detail I / II / III).
\end{frame}

% ---------- 7. Implementation ----------
\section{Implementation}
\begin{frame}{Implementation}
  \begin{itemize}
    \item Platform: simulator / RTL / compiler / runtime / hardware
    \item What is modeled vs. actually implemented; what is idealized
  \end{itemize}
\end{frame}

% ---------- 8. Evaluation Setup ----------
\section{Evaluation Setup}
\begin{frame}{Evaluation Setup}
  \begin{itemize}
    \item Baselines / workloads / metrics
    \item Hardware or simulation environment; fairness of comparison
  \end{itemize}
\end{frame}

% ---------- 9. Results ----------
\section{Results}
% Pattern for a hard result figure (split as needed):
%   Result Setup -> How to Read This Figure -> Main Trend -> Interpretation & Caveats
\begin{frame}{Result: <what is measured>}
  % \begin{figure}\includegraphics[width=0.85\linewidth]{figures/result1.png}\end{figure}
  % Label reused paper figures: "From the paper." / "Adapted from the paper."
  \begin{itemize}
    \item What question this answers; what is compared; axes
    \item Trend; authors' conclusion; whether it is justified; caveats
  \end{itemize}
\end{frame}

% ---------- 10. Strengths ----------
\section{Strengths \& Limitations}
\begin{frame}{Strengths}
  \begin{itemize}
    \item What the paper does well
  \end{itemize}
\end{frame}

% ---------- 11. Limitations ----------
\begin{frame}{Limitations}
  \begin{itemize}
    \item Assumptions / missing baselines / scalability / cost / mismatch
  \end{itemize}
\end{frame}

% ---------- 12. Discussion Questions ----------
\section{Discussion \& Takeaways}
\begin{frame}{Discussion Questions}
  \begin{enumerate}
    \item Question about assumptions / generality / methodology
    \item \dots (3--6 total; avoid trivially-answered questions)
  \end{enumerate}
\end{frame}

% ---------- 13. Takeaways ----------
\begin{frame}{Takeaways}
  \begin{itemize}
    \item Point 1
    \item Point 2
    \item Point 3
  \end{itemize}
\end{frame}

\begin{frame}{If You Remember Only One Thing\dots}
  \begin{center}
    \Large <the single most important takeaway>
  \end{center}
\end{frame}

% ---------- 14. Research Extensions ----------
\begin{frame}{Research Extensions}
  \begin{itemize}
    \item Concrete follow-up direction 1
    \item Concrete follow-up direction 2
  \end{itemize}
\end{frame}

\end{document}
